% !TeX root = main.tex
\chapter{Conclusion et perspectives}
\label{ch:conclusion}

Ce rapport a décrit le portage de StreamKM++ depuis l'implémentation Flink
de Bitsakis (2018) vers Apache Spark. Les trois briques algorithmiques de la
thèse --- Lloyd pondéré, seeding $k$-means++, coreset tree et
merge-and-reduce --- ont été implémentées et validées indépendamment de tout
moteur de calcul distribué (chapitre~\ref{ch:algorithmes}), avant d'être
orchestrées par une couche Spark qui ne contient elle-même aucune logique
algorithmique propre (chapitre~\ref{ch:implementation}). Cet ordre --- théorie,
puis framework, puis implémentation opérateur par opérateur, puis validation
de la qualité avant la scalabilité --- a été tenu tout au long du projet, et
a permis de découvrir et corriger un bug réel (une graine aléatoire mal
propagée dans la fusion de coresets) avant qu'il n'affecte les expériences
du chapitre~\ref{ch:evaluation}.

Le portage confirme ce que l'énoncé posait comme question implicite : une
partie du dataflow de la thèse se traduit directement (voire se simplifie,
comme la disparition de \texttt{HandleWatermark} ou la méthodologie
d'évaluation du coût), une partie se traduit avec un gain net (la fusion
arborescente des coresets, qui répond directement au goulot que la thèse
identifie elle-même), et une partie ne se traduit pas du tout (Queryable
State) --- sans que cela remette en cause la viabilité du portage dans son
ensemble (chapitre~\ref{ch:discussion}).

\section{Travaux futurs}

\paragraph{Variante à requêtes périodiques (§4.3 de la thèse).} Décrite mais
non implémentée faute de temps (chapitre~\ref{ch:implementation},
§\ref{sec:impl-requetes}) --- le chemin d'implémentation est cependant déjà
balisé : \texttt{Trigger.ProcessingTime} pour la ré-évaluation périodique,
\texttt{BucketManager.extractCoreset} (qui ne vide jamais l'état) pour
dupliquer le coreset courant à plusieurs valeurs de $k$.

\paragraph{Un vrai équivalent de Queryable State.} La seule perte de
fonctionnalité nette identifiée dans ce portage
(chapitre~\ref{ch:discussion}). Une piste non explorée : un service externe
léger interrogeant directement l'état exposé par \texttt{onModel}, plutôt
qu'un sink \texttt{memory} + SQL --- à évaluer en termes de fraîcheur et de
complexité ajoutée.

\paragraph{Gestion du concept drift.} Ajouter un mécanisme d'oubli
(fenêtre glissante sur les buckets les plus anciens, ou facteur de
décroissance à la \texttt{StreamingKMeans}) permettrait de traiter des flux
dont la distribution change réellement dans le temps --- hors du périmètre
de StreamKM++ tel que défini par la thèse, mais une extension naturelle.

\paragraph{Poursuivre l'optimisation du kernel local.} Le travail
Structure-of-Arrays / hors-tas engagé sur \texttt{core}/\texttt{coreset}
(chapitre~\ref{ch:evaluation}, benchmarks JMH) ouvre la voie à des
optimisations supplémentaires une fois les résultats consolidés --- en
particulier sur \texttt{CoresetTree.build}, identifié dès la spécification
initiale comme le point chaud principal du pipeline.
